\section{Moltiplicazione di matrici supercircolanti - Tr. di Hadamard}
\subsubsection*{Definizione di matrici supercircolanti}
Una matrice \(A \in \mathbb{C}^{n \times n}\), con \(n=2^k\) si dice supercircolante se:
\begin{itemize}
	\item per \(n = 1\) sempre (banalmente tutti i numeri sono considerati matrici supercircolanti)
	\item per \(n > 1\) la matrice si può dividere in quattro sottomatrici \(n/2 \times n/2\) uguali a due a due lungo la diagonale
	principale e la diagonale secondaria e tali sottomatrici sono a loro volta supercircolanti
	\[A = \begin{bmatrix} A_1 & A_2 \\ A_2 & A_1 \end{bmatrix} \]
\end{itemize}

\subsubsection*{Osservazioni sulle matrici supercircolanti}
\begin{itemize}
	\item nelle matrici supercircolanti, ogni riga è ottenuta tramite uno shift circolare della riga precedente
	\item le righe successive alla prima sono univocamente definite a partire dalla prima riga, detta riga generatrice,
	tramite opportuni shift circolari
	\item dato che è sufficiente la prima riga per definire l'intera matrice supercircolante, la taglia di una matrice
	supercircolante \(A\) è pari al numero di elementi della sua prima riga, ovvero \(n\)
	\item si osserva che l'insieme delle matrici supercircolanti è chiuso rispetto alla somma:
	\[A + B = \begin{bmatrix} A_1 & A_2 \\ A_2 & A_1 \end{bmatrix} + \begin{bmatrix} B_1 & B_2 \\ B_2 & B_1 \end{bmatrix} = \begin{bmatrix} A_1 + B_1 & A_2 + B_2 \\ A_2 + B_2 & A_1 + B_1 \end{bmatrix} \quad \text{supercircolante per def.}\]
	\item si osserva che l'insieme delle matrici supercircolanti è chiuso rispetto alla moltiplicazione:
	\[A \cdot B = \begin{bmatrix} A_1 & A_2 \\ A_2 & A_1 \end{bmatrix} \cdot \begin{bmatrix} B_1 & B_2 \\ B_2 & B_1 \end{bmatrix} = \begin{bmatrix} A_1 B_1 + A_2 B_2 & A_1 B_2 + A_2 B_1 \\ A_2 B_1 + A_1 B_2 & A_2 B_2 + A_1 B_1 \end{bmatrix} \quad \text{supercircolante per def.}\]
	\item siccome l'insieme delle matrici supercircolanti è chiuso rispetto alle operazioni di somma e moltiplicazione, allora
	è considerato un anello, in più si osserva che è anche commutativo, siccome le due operazioni godono della proprietà
	commutativa, ovvero \(A + B = B + A\) e \(A \cdot B = B \cdot A\)
\end{itemize}

\subsubsection*{Somme tra matrici supercircolanti}
Per sommare due matrici supercircolanti è sufficiente sommare le prime righe delle due matrici-addendi elemento per elemento, ottenendo
la prima riga della matrice risultante dalla somma. La complessità di questo approccio è \(T_{add}(n) = \Theta(n)\) con \(n\)
pari alla taglia delle matrici.

\subsection{Moltiplicazioni tra matrici supercircolanti}
\subsubsection*{Moltiplicazioni usando la definizione}
Usando la definizione di moltiplicazione di matrici, si osserva che si richiede il calcolo di 4 prodotti tra sottomatrici
supercircolanti di dimensione \(n/2 \times n/2\) e 2 somme tra matrici supercircolanti di dimensione \(n/2 \times n/2\).
L'equazione alle ricorrenze per questo approccio è la seguente:
\[T(n) \; = \; \begin{cases}
	\; 1 & \text{se } n = 1 \\[3pt]
	\; 4 \cdot T(n/2) + 2 \cdot n/2 & \text{se } n > 1
\end{cases} \;\; = \;\; \begin{cases}
	\; 1 & \text{se } n = 1 \\[3pt]
	\; 4 \cdot T(n/2) + n & \text{se } n > 1
\end{cases}\]
Dal master theorem con \(a = 4\), \(b = 2\) e \(\omega(n) = n\) si ottiene una funzione soglia di
\(\sigma(n) = n^{\log_2 4} = n^2\), tale per cui \(\omega(n) = \sigma(n) \cdot O(n^{-1})\), ovvero si ricade
nel primo caso del master theorem. La complessità è quindi \(T(n) = \Theta(\sigma(n)) = \Theta(n^2)\).

\newpage

\subsubsection*{Moltiplicazioni con approccio D\&C ottimizzato}
Si osserva che secondo opportune manipolazioni algebriche è possibile ricavare le due sottomatrici \(C_1\) e \(C_2\)
del prodotto attraverso soli 2 prodotti, con l'overhead di 6 somme e due moltiplicazioni per uno scalare.
\[\begin{array}{l}
	U = (A_1 + A_2) \cdot (B_1 + B_2) = A_1 B_1 + A_1 B_2 + A_2 B_1 + A_2 B_2 \\[3pt]
	V = (A_1 - A_2) \cdot (B_1 - B_2) = A_1 B_1 - A_1 B_2 - A_2 B_1 + A_2 B_2
\end{array} \qquad \begin{array}{l}
	C_1 = (U + V)/2 = A_1 B_1 + A_2 B_2 \\[3pt]
	C_2 = (U - V)/2 = A_1 B_2 + A_2 B_1
\end{array}\]
L'equazione alle ricorrenze per questo approccio è la seguente:
\[T(n) \; = \; \begin{cases}
	\; 1 & \text{se } n = 1 \\[3pt]
	\; 2 \cdot T(n/2) + 6 \cdot n/2 + 2 \cdot n/2 & \text{se } n > 1
\end{cases} \;\; = \;\; \begin{cases}
	\; 1 & \text{se } n = 1 \\[3pt]
	\; 2 \cdot T(n/2) + 4n & \text{se } n > 1
\end{cases}\]
Dal master theorem con \(a = 2\), \(b = 2\) e \(\omega(n) = 4n\) si ottiene una funzione soglia di
\(\sigma(n) = n^{\log_2 2} = n\) paragonabile a \(\omega(n)\). La complessità dal secondo caso del
master theorem risulta essere \(T(n) = \Theta(n \cdot \log n)\).

\subsection{Trasformata di Hadamard e Fast Hadamard Transform (FHT)}
\subsubsection*{Autovalori alle foglie dell'albero della ricorsione}
L'albero della ricorsione per il prodotto di matrici supercircolanti con approccio D\&C ottimizzato è un albero binario
completo in cui ogni nodo figlio sinistro riceve le somme tra le prime righe delle sottomatrici \(A_1\) e \(A_2\) e tra
quelle di \(B_1\) e \(B_2\), mentre ogni nodo figlio destro riceve le rispettive differenze, affinché siano moltiplicate opportunamente.

Le foglie ricevono due matrici degeneri in scalati \(a\) e \(b\) che è possibile dimostrare essere proprio gli autovalori
delle due matrici supercircolanti \(A\) e \(B\) di partenza.

\subsubsection*{Trasformata e prodotto di Hadamard}
È possibile interpretare le operazioni di divide come una trasformazione lineare che, partendo da una matrice supercircolante
rappresentata dalla sua prima riga, restituisce il vettore di autovalori della matrice supercircolante stessa. In questo modo
il prodotto di matrici supercircolanti può essere visto come un prodotto di autovalori, che risulta avere complessità
\(\Theta(n)\) con \(n\) pari alla taglia delle matrici supercircolanti (che è proprio il costo delle foglie). Una volta fatto
ciò, le operazioni di combine effettuano la trasformazione lineare inversa, restituendo la matrice supercircolante risultante
dal prodotto, rappresentata dalla sua prima riga.

La trasformazione lineare impiegata prende il nome di \textbf{trasformata di Hadamard}, mentre il prodotto membro a membro
tra i vettori di autovalori prende il nome di \textbf{prodotto di Hadamard}.

Come visto nell'analisi della complessità, la trasformata e la relativa antitrasformata di Hadamard vengono calcolate in
\(\Theta(n \cdot \log n)\), mentre il prodotto di Hadamard è calcolato in \(\Theta(n)\). Complessivamente, trasformata e
prodotto hanno complessità \(\Theta(n \cdot \log n)\) come l'approccio D\&C ottimizzato.

\subsubsection*{Cambio di rappresentazioni per risolvere problemi più efficientemente}
L'idea alla base del prodotto di matrici supercircolanti tramite la trasformata di Hadamard è quello di convertire i dati
in ingresso in una rappresentazione che permette di risolvere il problema richiesto in maniera più efficiente, per poi
convertire il risultato ottenuto nella rappresentazione originale. Altri esempi di algoritmi che adottano questa tecnica sono:
\begin{itemize}
	\item l'algoritmo di Strassen, anche se non è chiaro come interpretare la rappresentazione intermedia
	\item l'algoritmo di FFT, che utilizza la rappresentazione dei polinomi per punti ottenuta tramite la trasformata discreta
	di Fourier per calcolare il prodotto di polinomi in maniera più efficiente
\end{itemize}

\subsubsection*{Definizione di simmetria}
Una simmetria è una proprietà di invarianza rispetto ad una certa trasformazione. Ad esempio la proprietà commutativa è
l'invarianza rispetto alla trasformazione che scambia i due operandi di un'operazione.

Conoscere le simmetrie permette di identificare ed eliminare i gradi di libertà di un problema, riducendo sia il peso della
rappresentazione utilizzata che quello delle operazioni di calcolo.
